\uuid{5WE0}
\exo7id{1841}
\titre{exo7 1841}
\auteur{gourio}
\organisation{exo7}
\datecreate{2001-09-01}
\isIndication{false}
\isCorrection{false}
\chapitre{Extremum, extremum lié}
\sousChapitre{Extremum, extremum lié}
\cpge{mpsi-mp2i.fonctions-deux-variables.a, pcsi.fonctions-deux-variables.a, ptsi.fonctions-deux-variables.a, mpsi-mp2i.limites-continuite.c, pcsi.limites-continuite.c, ptsi.limites-continuite.c, mp-mpi.topologie-espaces-normes.h, mp-mpi.topologie-espaces-normes.j, pc.espaces-vectoriels-normes.g, psi.espaces-vectoriels-normes.g}
\module{Calcul différentiel}
\niveau{L3}
\difficulte{}

\contenu{
\question{
Soit $ f:{\Rr}^{2}\rightarrow {\Rr}$ continue telle que :
$$\lim\limits_{\left\| x\right\| \rightarrow \infty }\left| f(x)\right|
=+\infty . $$
Montrer que $ f$ est minor\'{e}e et atteint sa borne inf\'{e}rieure.
}
}
